PromptOps treats prompts as code: versioned, reviewed, and defended by a regression suite. The
prompt a model actually reads is not the prompt a team writes. It is first rewritten by the model's
\emph{chat template}, an executable Jinja program shipped inside the model repository, edited in
place, with no version of its own. How often this hidden dependency changes, and with what
behavioural consequence, has not been measured. We reconstruct the default-branch template history
of \NFrame{} of the most-downloaded text-generation repositories on the Hugging Face Hub. A
renderer validated byte-for-byte against the reference implementation decides whether each edit
changes what the model is shown, and we run \NCheckpoints{} local checkpoints under their own
repositories' historical templates. \PctDrifted\% of the repositories changed their template after
first publishing one, and \PctEventsWeightSilent\% of the change events moved no weight file,
invisible to a team that pinned the model and verified its weights; \PctEventsBeh\% of the
decidable edits change what the model is shown. Much of the raw rate is launch-week churn, and we
report the settled figure beside it: \PctBehSettledFrame\% of the frame saw such an edit land more
than \SettledDays{} days after its first template. In a one-repository case study, an
edit that dropped the system role halved output-contract compliance from \ConPhiA\% to \ConPhiB\%
until a later edit largely restored it. \textsc{templatelock} pins what a template renders rather than
what it says, separating the edits that matter from the ones that do not at \CostComputeMs{} ms per
model. Until prompt pipelines version the template, pinning the prompt and the model does not pin
the model's input.
